% Consolidated 27-optics-photometry from ff5921b6fbf473f9d16e47ad4b028ac5cf985609; original TeX fragments preserved.

% === BEGIN 27-optics-photometry-001-012.tex blob 935699ee1b34cd2db89bba05f3860bf976ee7947 ===
% Irodov Part V "Optics", §5.1, independently renumbered originals 1--12.
% Source uploaded PDF pp.196--198, printed pp.187--189. Figure 5.1.
\opticsproblem{1}{188}
参考原书图5.1的人眼相对光谱视见函数 $V(\lambda)$，
求：(a) 在 $\lambda=0.51,0.64\,\mathrm{\mu m}$ 两个单色波长下，产生 $1.0$lm 光通量所需的辐射能流；
(b) 波段 $0.58$—$0.63\,\mathrm{\mu m}$ 内辐射功率均匀分配，总辐射能流 $\Phi_e=4.5$mW，若 $V(\lambda)$ 在此段线性，求对应光通量。
\exsolution
原书采用峰值视见波长 $0.555\,\mathrm{\mu m}$ 的换算常量
$A=1.6\,\mathrm{mW/lm}$，
光通量与辐射能流满足
$\mathrm d\Phi_{\rm lum}=V(\lambda)\,\mathrm d\Phi_e/A$。
(a) 对单色光
$\Phi_e=A\Phi_{\rm lum}/V(\lambda)$。
按图5.1读数，
$V(0.51)\simeq0.53$、
$V(0.64)\simeq0.18$，有
\[
\boxed{\Phi_e(0.51)\simeq3\,\mathrm{mW}},\qquad
\boxed{\Phi_e(0.64)\simeq9\,\mathrm{mW}}.
\]
(b) 谱功率均匀、视见函数线性，故平均视见函数为端点算术平均：
\[
\boxed{\Phi_{\rm lum}
=\frac{\Phi_e}{2A}
[V(0.58)+V(0.63)]
\simeq1.6\,\mathrm{lm}}.
\]
\excheck 光度量与辐射能流不是同一单位；视见函数必须从原图读取，以上数值只在图上读数精度内有效。

\opticsproblem{2}{188}
频率对应波长 $\lambda=0.59\,\mathrm{\mu m}$ 的各向同性点光源，总光通量 $\Phi=10$lm。求距光源 $r=1.0$m 处电场与磁场强度的峰值，使用图5.1视见函数。
\exsolution
单色光总\emph{辐射}功率
$P_e=A\Phi/V(0.59)$，
图5.1该波长相对视见函数约为数十个百分点。
各向同性辐射在距离 $r$ 处的平均能流
$\overline S=P_e/(4\pi r^2)$。
真空平面波在局部远场满足
$\overline S=E_m^2/(2Z_0)$，
$Z_0=\sqrt{\mu_0/\varepsilon_0}$，
所以
\[
\boxed{E_m=
\sqrt{\frac{Z_0 A\Phi}{2\pi r^2V(0.59)}}
\simeq1.1\,\mathrm{V/m}},
\quad
\boxed{H_m=\frac{E_m}{Z_0}
\simeq2.9\,\mathrm{mA/m}}.
\]
\excheck 题中 $\Phi$ 为可见光\emph{光通量}而非总电磁波辐射功率，必须先借助 $V(\lambda)$ 换算。

\opticsproblem{3}{188}
求不透明球被照亮部分的平均照度：
(a) 受平行光照射，在法向入射点照度为 $E_0$；
(b) 各向同性点源距离球心 $\ell=100$cm，球半径 $R=60$cm，光源光强 $I=36$cd。
\exsolution
(a) 受照半球上纬度角为 $\theta$ 的面元照度 $E=E_0\cos\theta$。
半球平均
\[
\boxed{\bar E=\frac{\int_0^{2\pi}\!\!\int_0^{\pi/2}
E_0\cos\theta R^2\sin\theta\,\mathrm d\theta\mathrm d\varphi}
{2\pi R^2}=\frac{E_0}{2}}.
\]
(b) 点源看来球的张角满足
$\sin\alpha=R/\ell$，
被球遮挡的光通量
$\Phi=2\pi I(1-\cos\alpha)
=2\pi I[1-\sqrt{1-R^2/\ell^2}]$。
从球心看受照球冠面积
$S=2\pi R^2(1-R/\ell)$，因此
\[
\boxed{\bar E=
\frac{I[1-\sqrt{1-R^2/\ell^2}]}
{R^2(1-R/\ell)}
\simeq50\,\mathrm{lx}}.
\]
\excheck (b) 分母应是实际被照球冠面积，而不是半球面积；当光源远去时可恢复半球近似。

\opticsproblem{4}{188}
某发光面亮度随发射方向与法线夹角变化
$L(\theta)=L_0\cos\theta$，
求该面单位面积向上半空间发出的光出射度。
\exsolution
光出射度
$M=\int_{\Omega_+}L(\theta)\cos\theta\,\mathrm d\Omega$，
$\mathrm d\Omega=2\pi\sin\theta\,\mathrm d\theta$。
故
\[
M=2\pi L_0\int_0^{\pi/2}
\cos^2\theta\sin\theta\,\mathrm d\theta
=\boxed{\frac{2\pi}{3}L_0}.
\]
\excheck 方向亮度本身又带 $\cos\theta$，加上投影面积的 $\cos\theta$ 共形成平方，不能直接用朗伯源的 $\pi L_0$。

\opticsproblem{5}{189}
朗伯面光源的方向亮度为常数 $L$，求：
(a) 面元 $\Delta S$ 朝面法线、半顶角为 $\theta$ 的圆锥内的光通量；
(b) 光出射度 $M$。
\exsolution
面元向立体角微元发出
$\mathrm d\Phi=L\Delta S\cos\alpha\,\mathrm d\Omega$。
(a)
\[
\Phi_\theta=2\pi L\Delta S
\int_0^\theta\cos\alpha\sin\alpha\,\mathrm d\alpha,
\quad
\boxed{\Phi_\theta=\pi L\Delta S\sin^2\theta}.
\]
(b) 对上半球 $\theta=\pi/2$，
$\Phi=\pi L\Delta S$，所以
\[
\boxed{M=\Phi/\Delta S=\pi L}.
\]
\excheck 题中 $\theta$ 取\emph{圆锥半顶角}；若把开角理解为总顶角，须先除以2。

\opticsproblem{6}{189}
水平圆盘形朗伯光源面积
$S=100\,\mathrm{cm^2}$、亮度
$L=1.6\times10^4\,\mathrm{cd/m^2}$，
悬挂于半径 $R=1.0$m 的圆桌中心上方。
将圆盘视为点源，求使桌沿照度最大时的高度 $h$ 及最大照度。
\exsolution
从圆盘向桌沿的射线与竖直夹角 $\theta$，
$\cos\theta=h/\sqrt{R^2+h^2}$。
光源的方向光强 $I(\theta)=LS\cos\theta$，
桌面面元亦具有投影因子 $\cos\theta$，
所以
\[
E(h)=\frac{LS\cos^2\theta}{R^2+h^2}
=\frac{LSh^2}{(R^2+h^2)^2}.
\]
求导可得
$\mathrm dE/\mathrm dh\propto h(R^2-h^2)$，
故最大在 $h=R$，
\[
\boxed{h_*=R=1.0\,\mathrm m},\quad
\boxed{E_{\max}=\frac{LS}{4R^2}
=\frac{(1.6\times10^4)(0.010)}4
=40\,\mathrm{lx}}.
\]
\excheck 圆盘既是朗伯发射面又照亮水平桌面，方向光强和接受面各有一个 $\cos\theta$，不可仅保留一个。

\opticsproblem{7}{189}
点光源在半径 $R=1.0$m 的水平圆桌中心上方
$h=1.0$m。
希望桌面各点照度一致，求光强的方向函数
$I(\theta)$（$\theta$ 为出光方向与竖直的夹角），并当
$I(0)=I_0=100$cd 时求到达整个桌面的总光通量。
\exsolution
桌面距离为
$r=h/\cos\theta$，入射角亦为 $\theta$，
所以桌面照度
\[
E(\theta)=\frac{I(\theta)\cos\theta}{r^2}
=\frac{I(\theta)\cos^3\theta}{h^2}.
\]
要求 $E(\theta)=E(0)=I_0/h^2$，得
\[
\boxed{I(\theta)=\frac{I_0}{\cos^3\theta}},\qquad
0\le\theta\le\arctan(R/h).
\]
整个桌面光通量
\[
\boxed{\Phi=E(0)\pi R^2
=\frac{\pi I_0R^2}{h^2}
\simeq3.1\times10^2\,\mathrm{lm}}.
\]
\excheck 当 $\theta$ 增大，照射更远、更斜，为保持相同照度光强应\emph{增大}。书末文本似漏倒数，若照录 $I_0\cos^3\theta$ 则照度反而随 $\cos^6\theta$ 下降。

\opticsproblem{8}{189}
投影灯垂直照亮圆形房间天花板上的面积
$S=100\,\mathrm{cm^2}$ 的小斑点，斑点照度
$E=1000$lx，天花板朗伯反射率
$\rho=0.80$，房间半径 $R=2.0$m。
求反射斑点在圆柱形墙壁上造成的最大照度及其位置。
\exsolution
天花斑点向室内的光出射度
$M=\rho E$，对应朗伯亮度
$L=M/\pi=\rho E/\pi$。
在墙上、距天花垂直距离为 $z$ 的点，
与斑点的距离为 $d=\sqrt{R^2+z^2}$，
辐射源方向余弦 $z/d$，
墙法向余弦 $R/d$，故
\[
E_{\rm wall}(z)=\frac{\rho ES}{\pi}
\frac{Rz}{(R^2+z^2)^2}.
\]
微分得
$\mathrm dE/\mathrm dz\propto R^2-3z^2$，
最大位置
$\boxed{z_*=R/\sqrt3}$。
峰值
\[
\boxed{E_{\max}=
\frac{9}{16\sqrt3\pi}\frac{\rho ES}{R^2}
\simeq0.21\,\mathrm{lx}}.
\]
\excheck 书后纯文本位置标作 $R/3$，缺失根号；若用 $R/3$ 则不能使上述照度函数取最大。

\opticsproblem{9}{189}
一个半球形发光穹顶放在水平面上，表面光出射度处处相同，方向亮度为常数 $L$。求底面圆心处水平面元的照度。
\exsolution
从底面圆心看，穹顶占据完整上半球
$\Omega=2\pi$。
任一来光方向的辐亮度/亮度为 $L$，
水平接受面照度
\[
E=\int_{\Omega_+}L\cos\theta\,\mathrm d\Omega
=2\pi L\int_0^{\pi/2}
\cos\theta\sin\theta\,\mathrm d\theta.
\]
因此
\[
\boxed{E=\pi L}.
\]
\excheck 穹顶尺寸消去，因为每方向接收的亮度相同；假设穹顶发光均匀且无遮蔽。

\opticsproblem{10}{189}
亮度为常数 $L$ 的无限大朗伯平面，
求与该平面平行的任意小面元的照度。
\exsolution
任意接收点能见光源的完整上半空间，
且沿各入射方向的亮度为 $L$。
由照度的立体角积分
\[
\boxed{E=\int_0^{2\pi}\int_0^{\pi/2}
L\cos\theta\sin\theta\,
\mathrm d\theta\mathrm d\varphi=\pi L}.
\]
\excheck 对\emph{无限大}均匀光源，照度与两平面间距无关；有限发光平面无此性质。

\opticsproblem{11}{189}
均匀朗伯圆盘光源半径 $R=25$cm，
位于桌面上方 $h=75$cm。
圆盘轴下桌面照度 $E_0=70$lx，求光源光出射度 $M$。
\exsolution
圆盘均匀亮度为 $L=M/\pi$。
桌面轴上点对光源的立体角积分：
$E_0=2\pi L\int_0^R
[h^2r/(h^2+r^2)^2]\,\mathrm dr$，
积分得到
\[
E_0=\pi L\frac{R^2}{h^2+R^2}
=M\frac{R^2}{h^2+R^2}.
\]
因此
\[
\boxed{M=E_0\left(1+\frac{h^2}{R^2}\right)
=70(1+9)=700\,\mathrm{lm/m^2}}.
\]
\excheck 此处不能把圆盘当点源，原题提供的 $R/h$ 关系需在几何积分中保留。

\opticsproblem{12}{189}
半径 $R=6.0$cm 的均匀亮球灯，球心距地面
$h=3.0$m，亮度
$L=2.0\times10^4\,\mathrm{cd/m^2}$，
求球心正下方地面的照度。
\exsolution
小球在远处任一方向的表观投影面积为 $\pi R^2$，
其方向光强 $I=L\pi R^2$。
正下方的地面入射角为零，故
\[
\boxed{E_0=\frac{L\pi R^2}{h^2}}
=\frac{(2.0\times10^4)\pi(0.060)^2}{(3.0)^2}
\simeq25\,\mathrm{lx}.
\]
\excheck 当 $h\gg R$，球体可视为点源；计算采用\emph{投影圆盘面积}，不是球面总面积。
% === END 27-optics-photometry-001-012.tex ===

% === BEGIN 27-optics-photometry-013-025.tex blob 5357c4e5e84a118c42e8aa4b59afba2f95627094 ===
% Original Irodov Part V §5.1 problems 13--25, source printed pp.189--190.
\opticsproblem{13}{189}
用入射光线方向单位矢量 $\bm e$、反射光线方向 $\bm e'$ 及镜面外法线单位矢量 $\bm n$，写出反射定律的矢量形式。
\exsolution
分解入射矢量为切向与法向分量：
$\bm e=\bm e_t+(\bm e\cdot\bm n)\bm n$。
反射时切向保持、法向反号，故
\[
\boxed{\bm e'=\bm e-2(\bm e\cdot\bm n)\bm n}.
\]
\excheck $|\bm e'|=1$、$\bm e'\cdot\bm n=-\bm e\cdot\bm n$，分别验证速度方向模不变、入射与反射角相同。

\opticsproblem{14}{189}
证明依次在三个相互垂直的平面镜上反射后，出射光线方向恰与入射方向相反。
\exsolution
取三个镜法线 $\bm n_1,\bm n_2,\bm n_3$ 为正交单位基。
任一入射单位方向
$\bm e=\sum_{i=1}^3 e_i\bm n_i$，
在法线 $\bm n_i$ 的镜面反射会将对应分量 $e_i$ 变号，其他两个分量保持。
三次反射分别对三个分量变号，因此
\[
\boxed{\bm e_{\rm out}=-\bm e_{\rm in}}.
\]
\excheck 该结论不依赖入射方向，前提是光线依次确实击中三个镜面。

\opticsproblem{15}{189--190}
光从空气射向水面，求反射光线与折射光线互相垂直时的入射角。
\exsolution
两光线垂直意味着入射角 $\theta_B$ 与折射角 $r$ 满足
$\theta_B+r=90^\circ$。
斯涅尔定律
$\sin\theta_B=n\sin r=n\cos\theta_B$，
故
\[
\boxed{\theta_B=\arctan n\simeq53^\circ}
\]
（水折射率约 $n=1.33$）。
\excheck 正是布儒斯特角几何条件。

\opticsproblem{16}{190}
光从高折射率介质1射向介质2，
临界角为 $\theta_c$、反射与折射光垂直的入射角为 $\theta_B$，
已知 $\sin\theta_c/\sin\theta_B=\eta=1.28$，求 $n_1/n_2$。
\exsolution
记 $m=n_2/n_1<1$。全反射临界角
$\sin\theta_c=m$；布儒斯特角
$\tan\theta_B=m$，故
$\sin\theta_B=m/\sqrt{1+m^2}$。
于是 $\eta=\sqrt{1+m^2}$，得
\[
\boxed{\frac{n_1}{n_2}=\frac1{\sqrt{\eta^2-1}}
\simeq1.25}.
\]
\excheck 结果大于1，符合介质1光学更密；书末纯文本可能漏了根号。

\opticsproblem{17}{190}
光以入射角 $\theta=60^\circ$ 穿过厚
$d=6.0$cm 的平行玻璃板，求透射光的横向侧移。按普通玻璃折射率 $n=1.50$ 计算。
\exsolution
折射角 $r$ 满足 $\sin r=\sin\theta/n$。
侧移几何式
$\delta=d\sin(\theta-r)/\cos r$，
消元得
\[
\boxed{\delta=d\sin\theta
\left[1-\frac{\cos\theta}{\sqrt{n^2-\sin^2\theta}}\right]
\simeq3.1\,\mathrm{cm}}.
\]
\excheck 出射与入射光线平行；原题未单独给 $n$，数值以普通玻璃 $n\simeq1.5$ 为前提。

\opticsproblem{18}{190}
池底石头在水下真实深度为 $h$，
在空气中沿与法线夹角 $\theta$ 的方向观察。
求子午面内近轴光束形成的表观像深度。
\exsolution
令水折射率为 $n$，折射定律
$n\sin r=\sin\theta$。
主光线水面交点的横向位置
$x(\theta)=h\tan r=h\sin\theta/\sqrt{n^2-\sin^2\theta}$。
对邻近出射方向求导，
$\mathrm dx/\mathrm d\theta
=hn^2\cos\theta/(n^2-\sin^2\theta)^{3/2}$。
子午面小光束反向延长线交点深度
$h'=\cos^2\theta\,\mathrm dx/\mathrm d\theta$：
\[
\boxed{h'=\frac{hn^2\cos^3\theta}
{(n^2-\sin^2\theta)^{3/2}}}.
\]
\excheck $\theta\to0$ 时恢复 $h/n$；此处求的是小光束的子午像深度，而非单条光线的简单外推截距。

\opticsproblem{19}{190}
证明小顶角 $A$ 的棱镜在入射角亦很小时，偏折角一阶近似的模为 $(n-1)A$。
\exsolution
设第一面内外角 $i_1,r_1$，
第二面内外角 $r_2,i_2$。
几何 $r_1+r_2=A$，总偏折
$\delta=i_1+i_2-A$。
小角近似折射律给
$i_1\simeq nr_1$、
$i_2\simeq nr_2$，因而
\[
\boxed{|\delta|\simeq(n-1)A}.
\]
\excheck 原书取相反有向正角而写负号；
偏折大小和入射角的一阶无关性不受此选择影响。

\opticsproblem{20}{190}
证明棱镜对称穿过光路的偏折角最小，并给出最小偏折公式。
\exsolution
顶角 $A$，棱镜内两折射角满足
$r_1+r_2=A$；
外部两角
$i_j=\arcsin(n\sin r_j)$，
偏折 $\delta=i_1+i_2-A$。
在允许透射区间，
$f(r)=\arcsin(n\sin r)$ 的二阶导数
$f''(r)=n(n^2-1)\sin r/
(1-n^2\sin^2r)^{3/2}\ge0$。
凸性得
$f(r_1)+f(r_2)\ge2f(A/2)$，
因此最小在 $r_1=r_2=A/2$，
有 $i_1=i_2=(A+\delta_{\min})/2$：
\[
\boxed{\sin\frac{\delta_{\min}+A}{2}
=n\sin\frac A2}.
\]
\excheck 最小偏折对应入、出射角相等；必须存在两面都不全反射的实际传播光路。

\opticsproblem{21}{190}
某普通玻璃棱镜的最小偏折角恰等于顶角 $A$，求 $A$。
\exsolution
由第20题
$\sin[(\delta_{\min}+A)/2]=n\sin(A/2)$，
令 $\delta_{\min}=A$，
$\sin A=n\sin(A/2)$。
对于非零顶角，
$2\cos(A/2)=n$，所以
\[
\boxed{A=2\arccos(n/2)\simeq83^\circ}
\]
（按普通玻璃 $n\simeq1.50$）。
\excheck 原题未具体说明玻璃折射率，书末 $83^\circ$ 使用普通玻璃近似。

\opticsproblem{22}{190}
玻璃棱镜顶角 $A=60^\circ$，求偏折角的允许最小与最大值。采用普通玻璃 $n\simeq1.50$。
\exsolution
最小偏折：
\[
\delta_{\min}=2\arcsin[n\sin(A/2)]-A
=2\arcsin(0.75)-60^\circ\simeq37^\circ.
\]
在掠入射界限 $i_1\to90^\circ$ 时，
$r_1=\arcsin(1/n)$，
$r_2=A-r_1$，
$i_2=\arcsin(n\sin r_2)$，
因此
\[
\boxed{\delta_{\max}
=90^\circ+\arcsin[n\sin(A-\arcsin(1/n))]-A
\simeq58^\circ}.
\]
\excheck 最大值对应极限入射光线，实际表面透射功率会受到菲涅耳反射影响。

\opticsproblem{23}{190}
顶角 $60^\circ$ 的棱镜在空气中的最小偏折角为 $37^\circ$。把棱镜浸入水中，求新的最小偏折角。
\exsolution
空气中玻璃折射率
$n=\sin[(60^\circ+37^\circ)/2]/\sin30^\circ
\simeq1.50$。
水折射率取 $n_w\simeq1.33$；
浸水后的\emph{相对}折射率为 $n/n_w$，
所以
\[
\boxed{\delta_{w,\min}
=2\arcsin\left[\frac n{n_w}\sin30^\circ\right]
-60^\circ\simeq8.7^\circ}.
\]
\excheck 外介质改变必须将绝对折射率除以水折射率；直接照用空气中的 n 会高估偏折。

\opticsproblem{24}{190}
顶角 $A=60^\circ$ 的棱镜中两个单色分量折射率为 $n_1=1.515,n_2=1.520$。光路接近最小偏折，求出射角距。
\exsolution
最小偏折函数
$\delta(n)=2\arcsin[n\sin(A/2)]-A$，
微分：
$\mathrm d\delta/\mathrm dn=
2\sin(A/2)/\sqrt{1-n^2\sin^2(A/2)}$。
以平均折射率 $\bar n\simeq1.518$、
$\Delta n=0.005$ 计算：
\[
\boxed{\Delta\delta\simeq
\frac{2\sin(A/2)\Delta n}
{\sqrt{1-\bar n^2\sin^2(A/2)}}
\simeq0.0077\,\mathrm{rad}\simeq0.44^\circ}.
\]
\excheck 将弧度换成角度后才能与书末0.44°比较。

\opticsproblem{25}{190}
利用费马原理推导平面界面处光的反射和折射定律。
\exsolution
取界面上的可变光路交点横坐标为 $x$，
两段到源/目标的距离分别为 $r_1(x),r_2(x)$。
透射光程
$\mathcal L=n_1r_1+n_2r_2$，
费马原理令
$\mathrm d\mathcal L/\mathrm dx=0$。
几何微分
$\mathrm dr_1/\mathrm dx=\sin i$、
$\mathrm dr_2/\mathrm dx=-\sin r$，故
\[
\boxed{n_1\sin i=n_2\sin r}.
\]
若两段均在同一介质且终点位于反射侧，
同样变分得到
$n(\sin i-\sin i')=0$，
在通常角域内
$\boxed{i=i'}$。
\excheck 费马原理要求光程为\emph{驻值}，并非所有情形下严格的最短路径。
% === END 27-optics-photometry-013-025.tex ===

% === BEGIN 27-optics-photometry-026-035.tex blob c8055032c296615484869f5d44cfdb14b78d3aeb ===
% Irodov Part V §5.1 26--35. User PDF pp.199--201, printed 190--192,
% geometric figures 5.4--5.8 and source answers pp.325--326.
\opticsproblem{26}{190--191}
用作图法完成原书图5.4、5.5两组任务：
(a) 已知凹面和凸面球面镜、主轴及焦点，续作指定入射光线；
(b) 已知一对物像共轭点 P、P' 及主轴，确定镜面位置和焦点。
\exsolution
(a) 对球面镜使用两条基本射线：
与主轴平行的近轴入射光线反射后经过焦点（凸镜则其反向延长线经过虚焦点）；
指向曲率中心的光线沿原路返回。
从图5.4中已知焦点 $F$ 和镜顶点 $O$，曲率中心 $C$ 由
$\overline{OC}=2\overline{OF}$ 构造。
先作入射光线与镜面交点处的球面法线 $MC$，再将入射方向绕法线镜像，即得反射光线。
(b) 已知 P 与 P'，光源经镜面顶点反射的傍轴光线满足入射和反射角相等，可由共轭点的角平分线确定镜顶点；
再依有向球面镜公式
\[
\boxed{\frac1s+\frac1{s'}=\frac1f=\frac2R}
\]
确定焦距及曲率中心；对实、虚像用有向距离符号。
\excheck 原题为\emph{图解作图}，上述规则给出可执行的作图步骤，不捏造图5.4或5.5里未标出的数值。

\opticsproblem{27}{191}
凹面镜：(a) 物与像的距离 $l=15$cm、横向放大率
$\beta=-2.0$；
(b) 物体沿主轴移动 $l=5.0$cm 后，放大率由
$\beta_1=-0.50$ 变为 $\beta_2=-0.25$。
分别求镜焦距。
\exsolution
在实物实像约定下
$\beta=-s'/s$，
镜公式 $1/f=1/s+1/s'$。
(a) $s'=2s$，物像轴向间距
$|s'-s|=s=l$，
所以 $s=15$cm、$s'=30$cm，
\[
\boxed{f=\frac{ss'}{s+s'}=10\,\mathrm{cm}}.
\]
(b) 利用
$\beta=-f/(s-f)$，
可解物距
$s=f(1-1/\beta)$。
故两位置
$s_1=3f$、$s_2=5f$，
相距 $2f=l=5$cm，
\[
\boxed{f=2.5\,\mathrm{cm}}.
\]
\excheck 对实像放大率负号不可忽略；镜焦距为正的凹面镜光学绝对值。

\opticsproblem{28}{191}
点光源光强 $I_0=100$cd，距凹面镜顶点
$s=20.0$cm，焦距 $f=25.0$cm，镜面反射率
$\rho=0.80$。
求反射光束等效点像的光强。
\exsolution
球面镜成像公式
$1/s+1/s'=1/f$，
所以
$s'=fs/(s-f)
=(25)(20)/(20-25)=-100$cm，
形成放大率模
$|\beta|=|s'/s|=5$ 的虚像。
镜面反射后的同一束光功率是入射功率的 $\rho$ 倍，
从虚像发出的相同锥形方向固角缩小为原来的
$(s/|s'|)^2$ 倍（相当于等效辐射强度增加 $\beta^2$），所以
\[
\boxed{I'=\rho I_0
\left(\frac{f}{f-s}\right)^2
=0.80(100)(5)^2
=2.0\times10^3\,\mathrm{cd}}.
\]
\excheck 这里 $I'$ 为镜面反射束的\emph{等效像光强}，不是镜子新产生了额外总光通量。

\opticsproblem{29}{191}
从费马原理出发，推导折射率分别为
$n,n'$ 的两种介质之间的球面折射界面，对近轴光线满足的成像公式（界面曲率半径 $R$）。
\exsolution
将界面顶点置于原点，曲率球心在有向 $+x$ 方向为正，近轴表面高度 $h$ 处的轴向弓高
$x(h)\simeq h^2/(2R)$。
设物点、像点的有向轴距分别 $s,s'$，
按费马原理同一物像共轭光束的光程不随交点高度二阶改变。
从物点到界面、从界面到像点的光程展开至 $h^2$：
\[
n\sqrt{(x-s)^2+h^2}
+n'\sqrt{(s'-x)^2+h^2}
=\text{常数}+O(h^4).
\]
令 $h^2$ 系数为零，整理得到
\[
\boxed{\frac{n'}{s'}-\frac n{s}
=\frac{n'-n}{R}}.
\]
\excheck 公式的物像距离及曲率半径必须采用一致的有向正负约定；对 $n'=n$ 的界面，它退化为 $s'=s$ 而非产生折射光焦度。

\opticsproblem{30}{191}
图5.6：平行光从真空入射到折射率为 $n$ 的轴对称透明介质弯曲表面，欲使所有光线汇聚于顶点 O 后方轴向距离 $f$ 的 F 点。求界面形状 $x(r)$ 及可汇聚的最大入射束半径。
\exsolution
取 O 为 $(x,r)=(0,0)$，焦点 $F=(f,0)$。
与轴上入射光线相同光程的任一光线，先在真空中走到界面点 $(x,r)$，再在介质中到 F，满足
\[
x+n\sqrt{(f-x)^2+r^2}=nf.
\]
平方移项得到
$(n^2-1)x^2-2nf(n-1)x+n^2r^2=0$。
取 $x(0)=0$ 的根，
\[
\boxed{x(r)=\frac{nf}{n+1}
\left[1-\sqrt{1-\frac{(n+1)r^2}{(n-1)f^2}}\right]}.
\]
要求根号非负，因此
\[
\boxed{r_{\max}=f\sqrt{\frac{n-1}{n+1}}}.
\]
\excheck $r_{\max}$ 的长度量纲来自根号外的 $f$；书末答复的纯文本抽取对该根号不清晰，不宜误写为 $f(n-1)/(n+1)$。

\opticsproblem{31}{191}
点光源位于对称双凸玻璃厚透镜前表面
$20$cm 处；透镜厚 $d=5.0$cm，两面曲率半径的模均为
$R=5.0$cm。
求像位于透镜后表面之后多远。取普通玻璃
$n\simeq1.50$。
\exsolution
沿光传播方向取有向距离。
第一面从空气入玻璃，$R_1=+5$cm、物距 $s_1=-20$cm。
球面折射公式
$\frac{1.5}{s_1'}-\frac1{-20}
=\frac{0.5}{5}$（单位均取cm），
得到
$s_1'=30$cm。
第二面前的中间像位于其右方
$30-5=25$cm，故对第二面是虚物：
$s_2=+25$cm，
其曲率半径 $R_2=-5$cm，玻璃到空气：
\[
\frac1{s_2'}-\frac{1.5}{25}
=\frac{1-1.5}{-5}=0.10\,\mathrm{cm^{-1}},
\]
于是
$1/s_2'=0.16\,\mathrm{cm^{-1}}$，
\[
\boxed{s_2'=6.25\,\mathrm{cm}
\simeq6.3\,\mathrm{cm}}.
\]
\excheck 厚透镜必须分别处理前、后两个折射面，不能直接用薄透镜公式忽略5cm厚度。

\opticsproblem{32}{191--192}
图示厚 $d=9.0$cm 的平凸玻璃透镜，物体放在凸面前，
像恰形成于后方平面表面并以该面为屏。
凸面曲率半径 $R=2.5$cm，
物体亮度 $L=7700\,\mathrm{cd/m^2}$，
入瞳直径 $D=5.0$mm。
求：(a) 横向放大率；(b) 像照度。取玻璃折射率 $n=1.50$、无损耗。
\exsolution
设物体在空气中，进入玻璃后，像位于玻璃内距离前折射面 $d$ 处。
球面折射光学功率
$\Phi=(n-1)/R$，
后方像距 $s'=d$，
有向公式
$n/d-1/s=(n-1)/R$。
近轴横向放大率
$\beta=(n_{\rm in}s')/(n_{\rm out}s)
=d/(ns)$，
整理
\[
\boxed{\beta=1-\frac{d(n-1)}{nR}
=1-\frac{9(0.5)}{1.5(2.5)}
=-0.20}.
\]
以像方光束锥角和亮度守恒
$L_{\rm image}=n^2L$，
在小数值孔径近似下，
\[
\boxed{E\simeq\frac{\pi n^2D^2L}{4d^2}
\simeq42\,\mathrm{lx}}.
\]
\excheck 放大率负号表明实像倒立；亮度与面积比推导依赖近轴光束和忽略 Fresnel 损失。

\opticsproblem{33}{192}
求两种薄透镜的光焦度与物方、像方焦距：
(a) 玻璃透镜在空气中光焦度 $\Phi_0=-5.0$D，浸入
$n_0=1.7$ 的液体；
(b) 对称双凸玻璃薄透镜在空气中光焦度
$\Phi_0=+10$D，左侧为空气、右侧为水。
取玻璃折射率 $n\simeq1.50$、水折射率 $n_w\simeq1.33$。
\exsolution
(a) 相同面形的透镜光焦度与相对折射率差成正比：
\[
\Phi=\Phi_0\frac{n-n_0}{n-1}
=-5\frac{1.5-1.7}{1.5-1}=+2.0\,\mathrm D.
\]
均处于液体中，
$\boxed{f'=-f=n_0/\Phi=0.85\,\mathrm m}$。
(b) 两曲面光焦度之和
$\Phi=\frac{n-1}{R}+\frac{n_w-n}{-R}
=(2n-n_w-1)/R$；
空气中对称双凸 $\Phi_0=2(n-1)/R$，故
\[
\boxed{\Phi=\frac{\Phi_0(2n-n_w-1)}{2(n-1)}
\simeq6.7\,\mathrm D},
\]
$\boxed{f=-1/\Phi\simeq-15\,\mathrm{cm}}$，
$\boxed{f'=n_w/\Phi\simeq20\,\mathrm{cm}}$。
\excheck 双侧介质不同使物方与像方焦距绝对值不相同，不能照搬空气中 $f'=-f$。

\opticsproblem{34}{192}
按原书图5.7及图5.8作图：
(a) 会聚、发散薄透镜后的光线；
(b) 已知共轭点 P、P' 与主轴，求透镜位置和焦点；
(c) 已知某一条入射—出射光线对和主轴，求另一条入射光线的出射路径。
\exsolution
(a) 对会聚透镜，平行主轴入射光线出射后经过像方焦点 F'；
经过物方焦点 F 的入射光线出射后平行主轴；
通过光心 O 的近轴光线方向不变。
发散透镜应将“经过焦点”替换为“出射反向延长线经过虚焦点”。
(b) 连接给定共轭点 P、P'；这条不偏折的光心光线与主轴相交处即透镜光心，作与主轴垂直的透镜平面。
再用由 P 发出的平行主轴光线，结合 P' 连线与透镜平面交点求像方焦点 F'；倒作得 F。
(c) 将已知光线在透镜面上的交点及出射方向记录为主平面基准；
另一条入射光线可分解为穿过光心的不偏折光线与平行主轴光线，按共轭焦平面的同一点成像规则构造出射方向。
\excheck 本题是几何作图题而非数值题；给出了图5.7、5.8可操作的基准光线与焦平面构造规则，作图必须保持原图各射线位置关系。

\opticsproblem{35}{192}
焦距 $f=25$cm 的薄会聚透镜将物体像投影至距透镜 $l=5.0$m 的屏幕。屏幕向透镜方向移动
$\Delta l=18$cm。
问为使像重新清晰，物体应移动多少距离？
\exsolution
固定焦距，物距 $s$ 与像距 $l$ 满足
$1/f=1/s+1/l$，
即
$s(l)=fl/(l-f)$。
屏幕移至
$l'=l-\Delta l$，
物距新值
$s'=fl'/(l'-f)$，
因此精确位移为
\[
\Delta s=s'-s
=\frac{f^2\Delta l}{(l-f)(l-\Delta l-f)}.
\]
代 $f=0.25$m、$l=5.0$m、
$\Delta l=0.18$m，
\[
\boxed{\Delta s\simeq5.2\times10^{-4}\,\mathrm m
\approx0.5\,\mathrm{mm}}.
\]
\excheck 新像距变短后物距必须\emph{增加}，即物体远离透镜；近似 $\Delta s\simeq f^2\Delta l/(l-f)^2$。
% === END 27-optics-photometry-026-035.tex ===

% === BEGIN 27-optics-photometry-036-045.tex blob 2cbccf6475c540bfe1a363b63418fbf39f42d4d1 ===
% Original Irodov Part V §5.1 36--45, printed pp.193--194.
% Ray diagram 5.9 used for three-thin-lens system.
\opticsproblem{36}{193}
点光源与屏幕间固定距离 $l=90$cm，
一只薄会聚透镜在两个不同位置都可投得清晰实像。
求透镜焦距：
(a) 两个透镜位置间隔 $\Delta l=30$cm；
(b) 两位置得到的像横向大小之比 $\eta=4.0$。
\exsolution
两共轭透镜位置互换物距、像距：
$s+s'=l$，$|s'-s|=\Delta l$。
镜组透镜公式 $1/f=1/s+1/s'$ 给
$f=ss'/(s+s')$。
(a) $ss'=(l^2-\Delta l^2)/4$，所以
\[
\boxed{f=\frac{l^2-\Delta l^2}{4l}
=\frac{90^2-30^2}{4(90)}=20\,\mathrm{cm}}.
\]
(b) 两位置横向放大率模为
$\beta=s'/s$ 与其倒数；
比值 $\eta=\beta^2$，从而
$\beta=\sqrt\eta=2$。
取 $s=l/(1+\beta)=30$cm、
$s'=60$cm，得
\[
\boxed{f=\frac{ss'}l=20\,\mathrm{cm}}.
\]
\excheck 只在两个共轭实像位置存在时可用该方法，即 $l\ge4f$。

\opticsproblem{37}{193}
固定的物体与屏幕之间存在两种薄会聚透镜位置，
所得清晰像的横向尺寸分别为
$h_1'=2.0$mm 和 $h_2'=4.5$mm，
求物体本身横向尺寸 $h$。
\exsolution
两共轭透镜位置的横向放大率分别为
$\beta=s'/s$ 与
$\beta'=s/s'$，它们的乘积为1。
因此像大小之积
$h_1'h_2'=(\beta h)(h/\beta)=h^2$，故
\[
\boxed{h=\sqrt{h_1'h_2'}
=\sqrt{2.0\times4.5}\,\mathrm{mm}
=3.0\,\mathrm{mm}}.
\]
\excheck 物体大小是两个像尺寸的\emph{几何平均}，并非算术平均。

\opticsproblem{38}{193}
远处物体亮度 $L=260\,\mathrm{cd/m^2}$，
薄会聚摄影物镜孔径比
$D:f=1:3.5$；
透镜损耗光通量比例 $\alpha=0.10$。
求照片底片上成像照度。
\exsolution
对于远物，像距约为 $f$，透镜通光孔径在像点张的立体角近似
$\Omega\simeq\pi(D/2)^2/f^2$。
在近轴条件下像照度为
$E=(1-\alpha)L\Omega$，故
\[
\boxed{E=(1-\alpha)\frac{\pi LD^2}{4f^2}
=\frac{(0.9)\pi(260)}{4(3.5)^2}
\simeq15\,\mathrm{lx}}.
\]
\excheck 给的是摄影镜头的\emph{直径比} D/f，不是面积比；因此在照度公式中要平方。

\opticsproblem{39}{193}
薄会聚透镜成清晰实像，改变透镜通光口径 $D$，
讨论所成像的亮度：
(a) 眼睛直接观察真实空间像；
(b) 将像投到按朗伯定律漫反射的白屏上再观察。
\exsolution
(a) 在物镜光束完全进入眼瞳且未发生额外遮挡的条件下，
无耗损理想光学系统中辐亮度/亮度沿光线传播守恒。
改变口径主要改变通过的总光通量及视场，
并不改变被直接观察到的扩展实像表面亮度：
\[
\boxed{L_{\rm direct}\ \text{与}\ D\ \text{无关}}.
\]
(b) 屏上的照度由透镜收集的入射光通量决定，
通光面积 $\propto D^2$，
若其余条件固定，则
$E_{\rm screen}\propto D^2$。
朗伯白屏的反射亮度
$L_{\rm screen}=\rho E_{\rm screen}/\pi$，所以
\[
\boxed{L_{\rm screen}\propto D^2}.
\]
\excheck 两种观察方式须分别考虑“空间像的亮度不变量”与“漫反射屏上的照度”。

\opticsproblem{40}{193}
两只薄对称透镜紧密靠合，一只为双凸透镜
$n_1=1.70$，另一只为双凹透镜
$n_2=1.51$，两者表面曲率半径模均为
$R=10$cm。
将系统整体浸入水中，求像方焦距。
\exsolution
取水折射率 $n_w\simeq1.33$，
两透镜在水中的光焦度分别为
$\Phi_1=2(n_1-n_w)/R$，
$\Phi_2=-2(n_2-n_w)/R$。
紧贴时总光焦度
\[
\Phi=\Phi_1+\Phi_2=
\frac{2(n_1-n_2)}R
=\frac{2(1.70-1.51)}{0.10}
=3.8\,\mathrm D.
\]
介质中像方焦距 $f'=n_w/\Phi$，
因此
\[
\boxed{f'\simeq0.35\,\mathrm m=35\,\mathrm{cm}}.
\]
\excheck 水介质折射率在两透镜\emph{光焦度求和时}抵消，但求介质中的焦距还需乘 $n_w$。

\opticsproblem{41}{193}
取一只薄对称双凸玻璃透镜，
表面曲率半径模
$R=40$cm，将其中一面镀银，形成带折射面与反射面的凹面镜。
求等效焦距。
\exsolution
光依次在未镀银的凸面折射、
在银面反射、再从原折射面出射。
近轴薄系统的光焦度可按两次折射与一次反射的会聚功率相加，设玻璃折射率 $n\simeq1.50$：
\[
\Phi_{\rm eff}=
\frac{2(n-1)}R+\frac2R+
\frac{2(n-1)}R
=\frac{2(2n-1)}R.
\]
因而
\[
\boxed{f=\frac1{\Phi_{\rm eff}}
=\frac{R}{2(2n-1)}
=\frac{40\,\mathrm{cm}}{2(3-1)}
=10\,\mathrm{cm}}.
\]
\excheck 光两次穿过玻璃折射面，不能将其视为仅一面镀银的普通球面镜而用 $f=R/2$。

\opticsproblem{42}{193--194}
图5.9：三只共轴薄透镜依次排列，
光焦度为
$\Phi_1=+10.0$D、
$\Phi_2=-10.0$D、
$\Phi_3=+10.0$D；
相邻距离均为5.0cm。
求：(a) 从左平行射入的光束在最后一只透镜后汇聚的位置；
(b) 距第一透镜左侧多远有一点，与其像关于三透镜系统对称地位于两端？
\exsolution
使用光线列向量
$(y,u)^{\mathsf T}$，
自由传播和薄透镜矩阵分别为
$\mathsf P(d)=\begin{pmatrix}1&d\\0&1\end{pmatrix}$、
$\mathsf L(\Phi)=\begin{pmatrix}1&0\\-\Phi&1\end{pmatrix}$。
由左向右全系统
\[
\mathsf M=
\mathsf L(10)\mathsf P(0.05)
\mathsf L(-10)\mathsf P(0.05)\mathsf L(10)
=\begin{pmatrix}1/4&1/8\\-15/2&1/4\end{pmatrix}.
\]
(a) 平行光 $u_0=0$，
最后出射高度 $y'=y_0/4$、
斜率 $u'=-7.5y_0$，
向右传播到光轴需距离
\[
\boxed{s'=-y'/u'=\frac1{30}\,\mathrm m
\simeq3.3\,\mathrm{cm}}.
\]
(b) 若物点距第一透镜左侧 $\ell$，
像点距最后透镜右侧相同 $\ell$，
全传递矩阵
$\mathsf P(\ell)\mathsf M\mathsf P(\ell)$
的 (1,2) 元必须为零：
$-\tfrac{15}{2}\ell^2+\tfrac12\ell+\tfrac18=0$。
正根
\[
\boxed{\ell=\frac16\,\mathrm m\simeq17\,\mathrm{cm}}.
\]
\excheck 三透镜相隔各5cm，不能把总间距误当5cm；矩阵给出可重现的精确图示解。

\opticsproblem{43}{194}
伽利略望远镜角放大率为10，调焦到无穷远时筒长45cm。
求：(a) 物镜、目镜焦距；
(b) 观察50m远物体需将目镜移动多远。
\exsolution
伽利略望远镜物镜为正、目镜为负：
$\Gamma=f_{\rm obj}/|f_{\rm eye}|=10$，
调至无穷远筒长
$L=f_{\rm obj}-|f_{\rm eye}|=45$cm。
解得
\[
\boxed{f_{\rm obj}=50\,\mathrm{cm}},
\qquad
\boxed{f_{\rm eye}=-5.0\,\mathrm{cm}}.
\]
观察 $s=50$m 处物体，
物镜成像距
$s'=f_{\rm obj}s/(s-f_{\rm obj})$，
由原来0.500m变为
$0.50505$m，
因此
\[
\boxed{|\Delta x|=s'-f_{\rm obj}
\simeq0.00505\,\mathrm m=0.50\,\mathrm{cm}}.
\]
\excheck 远物变为有限距离时物镜中间像向右移动，需相应调整目镜的位置以恢复出射平行。

\opticsproblem{44}{194}
开普勒望远镜调至无穷远，物镜镜框直径为
$D$，经目镜成的镜框像直径为
$d$。求望远镜角放大率。
\exsolution
物镜所在平面经目镜的像称出瞳，
对于无穷远调整的望远镜，
光束的角放大率与光束口径缩小倍数互为对应：
$\Gamma=d\theta_{\rm out}/d\theta_{\rm in}
=f_{\rm obj}/f_{\rm eye}$。
从几何相似三角形得到出瞳直径
$d=D f_{\rm eye}/f_{\rm obj}$。
所以
\[
\boxed{\Gamma=\frac Dd}.
\]
\excheck 该结果需忽略遮挡及瞳孔截光，且把镜框的像正确视为出瞳。

\opticsproblem{45}{194}
望远镜使远方目标的光通量角集中度增大
$\eta=4.0\times10^4$ 倍。
目标经望远镜后的视角为
$\psi'=2.0^\circ$，
求目标自身视角 $\psi$。
\exsolution
在光损耗可忽略且面积源的小角度近似下，
角放大率为
$\Gamma=\sqrt\eta$。
因此
\[
\psi=\frac{\psi'}{\Gamma}
=\frac{2.0^\circ}{\sqrt{4.0\times10^4}}
=\frac{2^\circ}{200}
=0.01^\circ.
\]
\[
\boxed{\psi=0.60'}
\]
其中角分 $1^\circ=60'$。
\excheck 不能把能量强度倍数4万直接当角放大率；角尺度对应平方根关系。
% === END 27-optics-photometry-036-045.tex ===

% === BEGIN 27-optics-photometry-046-057.tex blob 226e878f671b8abe73956cc84f21a0631b8805a8 ===
% Irodov Part V §5.1 original 46--57. Source printed pp.194--196,
% original diagrams Fig.5.10--5.11; source answers printed 326--327.
\opticsproblem{46}{194}
开普勒望远镜原角放大率 $\Gamma=15$。
望远镜内外充满水，但筒长不变，为重新实现望远镜无穷远调整，只更换物镜。
目镜玻璃折射率为 $n=1.50$。求新角放大率。
\exsolution
原目镜在空气中焦距为 $f_e$，
原物镜 $f_o=\Gamma f_e$，
筒长 $\ell=(\Gamma+1)f_e$。
浸入折射率 $n_w\simeq1.33$ 的水后，
同一玻璃薄目镜焦距
$f_{e,w}=f_e\,n_w(n-1)/(n-n_w)$，
记其倍率为
$K=n_w(n-1)/(n-n_w)$。
新物镜调整为
$f_{o,w}=\ell-f_{e,w}
=(\Gamma+1-K)f_e$。
故
\[
\boxed{\Gamma_w=
\frac{f_{o,w}}{f_{e,w}}
=(\Gamma+1)\frac{n-n_w}{n_w(n-1)}-1
\simeq3.1}.
\]
\excheck 系统筒长保持不变意味着\emph{更换物镜}，不是简单将两焦距按同一个比例缩放。

\opticsproblem{47}{194}
望远镜物镜口径
$D=6.0$cm，
人眼瞳孔直径
$d_0=3.0$mm。
忽略光损耗，问角放大率 $\Gamma$ 最多可以多大，
以使视网膜像照度不低于裸眼观察？
\exsolution
由望远镜光束的拉格朗日不变量，
物镜通光口径 $D$ 在目镜后形成的出瞳直径
$d_{\rm exit}=D/\Gamma$。
若出瞳比眼瞳小，则有部分裸眼可接受光锥面积被压缩，
会降低扩展目标的视网膜照度。
不降低照度的条件
$d_{\rm exit}\ge d_0$，所以
\[
\boxed{\Gamma\le\frac D{d_0}
=\frac{60\,\mathrm{mm}}{3.0\,\mathrm{mm}}
=20}.
\]
\excheck 该约束适用于扩展目标、忽略损耗与无杂散光；不能不加限制套用于单点星等探测。

\opticsproblem{48}{194}
显微镜物镜和目镜的光焦度分别是
$100$D、$20$D，
当前角放大率 $\Gamma=50$。
两镜间距增加 $2.0$cm，问新放大率。
\exsolution
物镜焦距 $f_o=1/100$ m $=1.0$cm，
目镜焦距 $f_e=1/20$m $=5.0$cm。
显微镜近轴放大率
$\Gamma\simeq\Delta\,\ell_0/(f_of_e)$，
其中 $\Delta$ 为光学管长，
$\ell_0=25$cm 为标准明视距离。
间距增加 $\delta\Delta=2.0$cm，
增量
\[
\delta\Gamma\simeq\frac{\ell_0\delta\Delta}
{f_of_e}
=\frac{25(2.0)}{(1.0)(5.0)}
=10.
\]
故
\[
\boxed{\Gamma_{\rm new}\simeq60}.
\]
\excheck 这里保持物镜、目镜焦距不变，假设中间实像依然位于目镜可调整的成像条件下。

\opticsproblem{49}{194}
显微镜物镜数值孔径
$\sin\alpha=0.12$，
眼瞳直径
$d_0=4.0$mm。
确定放大率
(a) 使进入眼瞳的光束直径恰等于眼瞳直径；
(b) 使视网膜像照度不依赖显微镜放大率（受物镜框限光时）。
\exsolution
小孔径近轴模型中，经过显微镜后到达眼瞳的光束直径约
$d_{\rm exit}\simeq2\alpha\ell_0/\Gamma$，
$\ell_0=25$cm 是明视距离。
(a) 令
$d_{\rm exit}=d_0$，
有
\[
\boxed{\Gamma_*=\frac{2\alpha\ell_0}{d_0}
\simeq\frac{2(0.12)(0.25)}{0.004}
=15}.
\]
(b) 当 $\Gamma\le\Gamma_*$，
出射光束直径不小于眼瞳，瞳孔限光使视网膜像亮度近似不再受放大率影响，
\[
\boxed{\Gamma\le15}.
\]
\excheck 若 $\Gamma>15$，显微镜出瞳小于眼瞳，改变放大率会改变视网膜接收的照度与角面积关系。

\opticsproblem{50}{194}
一只薄对称双凸玻璃透镜两面曲率半径模
$R=7.50$cm，
一侧为空气、一侧为水。
求两个主平面、物方及像方焦点、节点位置。
取玻璃折射率 $n\simeq1.50$、水折射率
$n_w\simeq1.33$。
\exsolution
薄透镜的两个主平面重合于其中心面 H。
空气入射、水出射时总光焦度
\[
\Phi=\frac{n-1}R+\frac{n_w-n}{-R}
=\frac{2n-n_w-1}{R}
\simeq\frac{0.67}{0.075}
\simeq8.9\,\mathrm D.
\]
物方、像方焦距分别为
\[
\boxed{f=-\frac1\Phi\simeq-11\,\mathrm{cm}},
\quad
\boxed{f'=\frac{n_w}{\Phi}\simeq+15\,\mathrm{cm}}.
\]
入射、出射介质折射率不同，两个节点不再与主点重合；
二节点重合于由光学节点法则确定的共轴位置，
其偏离透镜平面的有向量为
\[
\boxed{x_N=f'+f\simeq3.7\,\mathrm{cm}}
\]
（在水侧）。
\excheck 薄透镜两\emph{主平面}重合不意味着两个\emph{节点}也恰在透镜中心；需考虑两侧折射率差。

\opticsproblem{51}{195}
原书图5.10给出三个共轴系统，
按图示作出焦点和主平面：
(a) 由焦距 $f_1=1.5a$、
$f_2=-1.5a$ 的会聚与发散薄透镜组成的远摄系统，间距为 $a$；
(b) 由两只会聚薄透镜
$f_1=1.5a$、
$f_2=0.5a$ 组成，间距仍为 $a$；
(c) 厚凹凸透镜两面光焦度为
$\Phi_1=+50$D、
$\Phi_2=-50$D，
两面间距 $d=4$cm、内部折射率 $n=1.5$。
\exsolution
对相距 $d$、中间介质折射率 $n$ 的两折射面，
总光焦度
$\Phi=\Phi_1+\Phi_2-(d/n)\Phi_1\Phi_2$；
主平面相对前后面的有向距离
\[
\boxed{x=\frac{d\Phi_2}{n\Phi}},\qquad
\boxed{x'=-\frac{d\Phi_1}{n\Phi}}.
\]
(a) 空气中
$\Phi_1=2/(3a)$、
$\Phi_2=-2/(3a)$，
$\Phi=4/(9a)$，
所以
$\boxed{f=9a/4}$，
$\boxed{x=-3a/2,\ x'=-3a/2}$。
(b)
$\Phi_1=2/(3a)$、
$\Phi_2=2/a$，
$\Phi=4/(3a)$，
于是
$\boxed{f=3a/4}$，
$\boxed{x=3a/2,\ x'=-a/2}$。
(c)
$\Phi=-(d/n)(+50)(-50)=66.7$D，
$\boxed{f=1/\Phi=1.5\,\mathrm{cm}}$，
$\boxed{x=-2.0\,\mathrm{cm},\ x'=-2.0\,\mathrm{cm}}$。
依这些有向位置分别在原图轴线上竖起主平面，
并从 H、H' 向物像两侧量取焦距即可。
\excheck 主平面可落在实体镜组\emph{外部}；其有向位置不能仅凭镜组外形猜测。

\opticsproblem{52}{195}
按图5.11(a)(b)(c)所给光轴、焦点、主平面、部分已知共轭点或一条入射—出射光线，用几何作图求缺失的
(a) 像方焦点 F' 与主平面 H'；
(b) 点 S 的像点 S'；
(c) 前后焦点 F、F' 与后主平面 H'。
\exsolution
(a) 由已知平行轴入射光线确定主平面上等高对应点，
其折射后的延长线与轴线交于像方焦点 F'；
在给定任一非轴光线的入、出射延长线上截取相同高度点，分别作垂直轴线的前后主平面。
(b) 由 S 发出一条平行轴的辅助光线，在前主平面 H 上标记高度；
在后主平面 H' 相同高度处连接 F' 得其出射线；
再取由 S 指向前焦点 F 的第二辅助光线，其经系统后平行主轴出射；
两条出射线（必要时用反向延长线）的交点为 S'。
(c) 已知图中一对入射—出射光线的延长线，
在光轴方向上等高对应构造主平面；
平行轴入射光在后焦平面处汇聚定义 F'，
反向追迹平行轴出射光确定 F，
从而补齐图示标注。
\excheck 对图解题，解法的可检验对象是两条主光线的几何共轭条件；任何具体位置都须依据原图给定相对坐标作图，本稿不伪造未给出的长度。

\opticsproblem{53}{195}
按原书列出的四种共轴顺序
(a) F—S—H—H'—F'；
(b) H—S—F'—F—H'；
(c) H'—S—F'—F—H；
(d) F'—H'—S—H—F，
分别作出物点 S 的像点 S'。
\exsolution
统一利用物像共轭的两主光线：
第一条从 S 平行光轴到前主平面 H，
在后主平面 H' 的同一高度处转为指向 F' 的出射线；
第二条从 S 沿指向前焦点 F 的方向入射，
在后主平面 H' 同一高度点转为平行光轴的出射线。
两出射线（或其反向延长线）的交点定义像点 S'。
任一情况下都可用有向牛顿共轭式代数核对：
\[
\boxed{xx'=ff'},
\]
其中 $x$ 是 S 到物方焦点 F 的有向距离，
$x'$ 为像点 S' 到像方焦点 F' 的有向距离。
(a) S 位于 F 与 H 间，出射线反向延长线相交，得到物方虚像；
(b)—(d) 按相应给定点序，从 F 量取 $x$，利用 $x'=ff'/x$ 在 F' 的正确有向一侧标出 S'，再以两主光线交点复核。
\excheck 这四组符号顺序涉及非通常次序的主平面和焦点，必须使用有向距离；不能简单地把所有光路解释成普通薄会聚透镜。

\opticsproblem{54}{195}
远摄镜组由前会聚透镜
$\Phi_1=+10$D 与后发散透镜
$\Phi_2=-10$D 组成。
求：(a) 间距 $d=4.0$cm 时的总焦距、主平面位置；
(b) 使总焦距 $f$ 与前透镜至后焦点距离
$\ell$ 的比值最大时的间距及最大比值。
\exsolution
镜组位于空气中，
总光焦度
$\Phi=\Phi_1+\Phi_2-d\Phi_1\Phi_2=100d$
（$d$ 用米），故
$\boxed{f=1/(100d)}$。
(a) $d=0.04$m 时
$\Phi=4$D，
$\boxed{f=25\,\mathrm{cm}}$。
主平面位置（相对对应透镜顶点）：
\[
\boxed{x=\frac{d\Phi_2}\Phi=-10\,\mathrm{cm}},
\qquad
\boxed{x'=-\frac{d\Phi_1}\Phi=-10\,\mathrm{cm}}.
\]
(b) 后主平面距后透镜有向距离 $-0.10$m；
后焦点在从前透镜向右 $d-0.10+f$ 处，
所以
$\ell=d-0.10+1/(100d)$。
比值
\[
\frac f\ell=\frac1{100d^2-10d+1}.
\]
分母在 $d=0.05$m 最小，故
\[
\boxed{d_{\max}=5.0\,\mathrm{cm}},
\quad
\boxed{(f/\ell)_{\max}=\frac43}.
\]
\excheck 远摄镜组的有效焦距可以大于从前镜到焦点的机械长度；这正是远摄压缩结构的目的。

\opticsproblem{55}{195--196}
厚弯月形玻璃透镜的凸面曲率半径
$R_1=10.0$cm、凹面
$R_2=5.0$cm，厚
$d=3.0$cm。
求前后主平面和前、后焦点的位置。取玻璃 $n=1.50$、外为空气，并规定光从凸面入射。
\exsolution
前面有向曲率半径 $R_1=+0.10$m，
后面凹面按光从左向右传播的曲率中心同在右侧，
$R_2=+0.05$m。
两面光焦度
$\Phi_1=(n-1)/R_1=+5$D、
$\Phi_2=(1-n)/R_2=-10$D。
厚透镜总焦度
\[
\Phi=\Phi_1+\Phi_2-\frac d n\Phi_1\Phi_2
=5-10-\frac{0.03}{1.5}(5)(-10)
=\boxed{-4.0\,\mathrm D}.
\]
因此有效像方焦距
$\boxed{f'=-25\,\mathrm{cm}}$，
物方焦距 $f=+25$cm。
两主平面相对于凸面顶点和凹面顶点的有向位置为
\[
\boxed{x=\frac{d\Phi_2}{n\Phi}=+5.0\,\mathrm{cm}},
\quad
\boxed{x'=-\frac{d\Phi_1}{n\Phi}=+2.5\,\mathrm{cm}}.
\]
按上述有向约定先标主平面，随后在各主平面向物方、像方量取有向焦距即可确定焦点。
\excheck 弯月透镜总体为发散型；
主平面并不必处于厚透镜实体内部，前后面的法线和曲率号必须保持一致。

\opticsproblem{56}{196}
两只共轴薄透镜焦距为
$f_1,f_2$，相距 $d$。
希望用单只薄透镜代替整个镜组，
使任意物距处的横向放大率与原系统相同。
求替代透镜焦距与其相对第一只镜的位置。
\exsolution
两薄透镜在空气中的光焦度
$\Phi=1/f_1+1/f_2-d/(f_1f_2)$。
等效焦距
\[
\boxed{f_{\rm eq}=\frac{f_1f_2}{f_1+f_2-d}}.
\]
要在任意物距保持横向放大率，不能只匹配光焦度；
还须把替代透镜的主平面置于原镜组的前主平面 H。
该主平面距第一只镜的有向距离
\[
\boxed{x=\frac{d(1/f_2)}{\Phi}
=\frac{f_1d}{f_1+f_2-d}}.
\]
\excheck 若 $f_1+f_2-d=0$，镜组为无焦系统，
替代焦距趋于无穷，此时无法用普通有限焦距单薄透镜描述。

\opticsproblem{57}{196}
一只薄对称双凸玻璃透镜表面曲率半径模
$R=38$cm；
垂直光轴的平面镜在透镜后方
$\ell=12$cm 处，
透镜与镜间充满水。
求光从空气侧入射、经水区镜面反射再返回的等效光焦度。
\exsolution
取玻璃折射率 $n=1.50$、水折射率
$n_w\simeq1.33$。
从空气入射、经透镜进入水中时单次透镜光焦度
\[
\Phi'=\frac{n-1}{R}+
\frac{n_w-n}{-R}
=\frac{2n-n_w-1}{R}.
\]
光往返两次经过透镜，并穿越水层至平面镜返回，
其等效反射光焦度是
\[
\boxed{\Phi_{\rm eff}
=2\Phi'-\frac{2\ell}{n_w}(\Phi')^2}.
\]
代 $R=0.38$m、$\ell=0.12$m，
$\Phi'=(3-1.33-1)/0.38\simeq1.76$D，
得到
\[
\boxed{\Phi_{\rm eff}\simeq
2(1.76)-\frac{0.24}{1.33}(1.76)^2
\simeq3.0\,\mathrm D}.
\]
\excheck 平面镜本身不提供曲率光焦度，
但光往返水层的传播距离进入等效光焦度的二次修正项。
% === END 27-optics-photometry-046-057.tex ===

% === BEGIN 27-optics-photometry-058-063.tex blob 6144e92f5f0636fb3667b7bad79b2553d4e6bef1 ===
% I. E. Irodov, Problems in General Physics, Part V §5.1, original 58--63.
% Verified source: uploaded PDF pages 205--206, printed 196--197; book answers PDF ~336.
\opticsproblem{58}{196}
一片空气中的厚玻璃弯月透镜，凸面曲率半径为 $R_1$、凹面为 $R_2$，凸面半径比凹面大 $\Delta R=1.5\,\mathrm{cm}$。(a) 求成为无焦望远系统时所需厚度。(b) 若两曲率半径分别为 $10.0$ 与 $7.5\,\mathrm{cm}$，求总光焦度为 $-1.0\,\mathrm D$ 时的厚度。
\exsolution
以入射方向为正，厚度为 $d$，玻璃折射率取 $n\simeq1.50$。两面的面光焦度为
$\Phi_1=(n-1)/R_1$、
$\Phi_2=-(n-1)/R_2$。厚透镜总光焦度
\[
\Phi=\Phi_1+\Phi_2-\frac dn\Phi_1\Phi_2
=(n-1)\left(\frac1{R_1}-\frac1{R_2}\right)
+\frac{d(n-1)^2}{nR_1R_2}.
\]
(a) 无焦条件 $\Phi=0$：
\[
\boxed{d=\frac{n(R_1-R_2)}{n-1}
=\frac{n\Delta R}{n-1}=4.5\,\mathrm{cm}}.
\]
(b) 由上式解 $d$，
\[
d=\frac{nR_1R_2}{(n-1)^2}
\left[\Phi-(n-1)\left(\frac1{R_1}-\frac1{R_2}\right)\right],
\]
代入 $R_1=0.100$m、$R_2=0.075$m、$\Phi=-1\,\mathrm D$，
\[
\boxed{d=3.0\,\mathrm{cm}}.
\]
\excheck 厚度单位在光焦度方程中必须为米；令 $d=0$ 还原为普通发散弯月薄透镜。

\opticsproblem{59}{196}
求厚玻璃弯月透镜的主平面位置、焦距及光焦度的正负：(a) 厚度 $d$，两面的曲率半径模均为 $R$；(b) 两面同心，半径分别为 $R_1,R_2$，且 $R_2>R_1$。
\exsolution
总光焦度与主平面位置满足
\[
\Phi=\Phi_1+\Phi_2-\frac dn\Phi_1\Phi_2,\quad
x=\frac{d\Phi_2}{n\Phi},\quad
x'=-\frac{d\Phi_1}{n\Phi},
\]
$x,x'$ 分别从前、后表面顶点沿入射方向量取。
(a) $\Phi_1=(n-1)/R$、
$\Phi_2=-(n-1)/R$，得
\[
\boxed{\Phi=\frac{d(n-1)^2}{nR^2}>0},\qquad
\boxed{f'=\frac1\Phi=\frac{nR^2}{d(n-1)^2}},
\]
\[
\boxed{x=x'=-\frac R{n-1}}.
\]
因此前后主平面相距 $d$，都位于凸面一侧的外部。
(b) 同心两表面的几何曲率中心相同。以第一折射面曲率中心及第二面半径的几何符号代入，消去 $d=R_2-R_1$：
\[
\boxed{\Phi=\frac{n-1}{n}
\left(\frac1{R_2}-\frac1{R_1}\right)<0}.
\]
两个主平面通过共同曲率中心，系统为发散型，有向焦距 $\boxed{f'=1/\Phi<0}$。
\excheck (a) 中虽然两面的薄透镜光焦度和为零，有限厚度的传播项仍形成正光焦度；(b) 中 $R_2>R_1$ 确保所给光焦度为负。

\opticsproblem{60}{196}
由半径 $R_1=5.0$cm、$R_2=1.0$cm 的两个均匀玻璃球透镜组成无焦望远系统。较大球作物镜，求两球中心距离与角放大率。
\exsolution
折射率为 $n$ 的玻璃球在空气中可作球透镜，其等效焦距（相对球心主平面）为
$f=nR/[2(n-1)]$。两球球心即各自对称主平面位置。
无焦望远条件使两球间距
$d=f_1+f_2$：
\[
\boxed{d=\frac{n(R_1+R_2)}{2(n-1)}
\simeq9.0\,\mathrm{cm}}\quad(n=1.50).
\]
系统角放大率的模
\[
\boxed{|\Gamma|=\frac{f_1}{f_2}=\frac{R_1}{R_2}=5.0}.
\]
\excheck 两球均为会聚元件，对无穷远望远系统有倒像；若要保留方向，带符号角放大率为 $-5$。

\opticsproblem{61}{196}
两只完全相同的厚对称双凸玻璃透镜紧密相贴。每只透镜厚度等于球面曲率半径 $d=R=3.0$cm，求组合系统在空气中的光焦度。
\exsolution
每个球面折射的光焦度模为
$\phi=(n-1)/R$。以约化角为变量的光线矩阵分别为
\[
\mathsf S(\phi)=\begin{pmatrix}1&0\\-\phi&1\end{pmatrix},\quad
\mathsf T=\begin{pmatrix}1&R/n\\0&1\end{pmatrix}.
\]
单片透镜的矩阵是
$\mathsf M=\mathsf S(\phi)\mathsf T\mathsf S(\phi)$，
组合后是 $\mathsf M^2$。
取系统矩阵下左元的负数作为空气中等效光焦度，并令 $d=R$，代数整理：
\[
\boxed{\Phi_{\rm tot}=\frac{2(n^2-1)}{n^2R}}.
\]
以 $n=1.50$、$R=0.030$m 算得
\[
\boxed{\Phi_{\rm tot}\simeq37\,\mathrm D}.
\]
\excheck 若两厚透镜中心厚度不可忽略，不能只把两个零厚薄透镜光焦度简单相加；透镜矩阵保留了内部传播效应。

\opticsproblem{62}{196}
各向同性但折射率 $n(\bm r)$ 缓慢变化的介质中，一条光线曲率半径 $\rho$ 满足
$1/\rho=\partial_N(\ln n)$；其中 $\partial_N$ 为沿光线主法线的方向导数。利用局部折射定律推导。
\exsolution
沿光线弧长记单位切线 $\bm t=\mathrm d\bm r/\mathrm ds$，费马原理在光学非均匀介质中的光线方程为
\[
\frac{\mathrm d}{\mathrm ds}(n\bm t)=\nabla n.
\]
展开左端：
$(\mathrm dn/\mathrm ds)\bm t+
n\,\mathrm d\bm t/\mathrm ds=\nabla n$。
取垂直 $\bm t$ 的主法线 $\bm N$ 的投影，第一项消失。
弗雷内公式 $\mathrm d\bm t/\mathrm ds=\bm N/\rho$，
从而
\[
\frac n\rho=\nabla n\cdot\bm N=\partial_N n,
\qquad
\boxed{\frac1\rho=\partial_N(\ln n)}.
\]
\excheck 曲率朝折射率增大的方向，若折射率梯度垂直于光线为零，则曲率为零，恢复直线传播。

\opticsproblem{63}{196--197}
地表附近水平方向传播的光线，空气折射率沿高度的梯度模为
$|\nabla n|\simeq3\times10^{-8}\,\mathrm{m^{-1}}$。
求光线的曲率半径；若光线要沿地球表面弯曲绕地一周，需多大的折射率梯度？
\exsolution
由上题
$\rho^{-1}=|\partial_N\ln n|\simeq|\nabla n|/n$。
空气 $n\simeq1$，故
\[
\boxed{\rho\simeq\frac1{3\times10^{-8}}
\simeq3.3\times10^7\,\mathrm m}.
\]
若光线与地球表面保持同曲率，
$\rho=R_\oplus\simeq6.4\times10^6$m，
要求
\[
\boxed{|\nabla n|_{\rm req}\simeq
\frac1{R_\oplus}\simeq1.6\times10^{-7}\,\mathrm{m^{-1}}}.
\]
\excheck 需折射率朝地面方向增大才可使光线朝地心弯曲；只给梯度大小不确定方向。
% === END 27-optics-photometry-058-063.tex ===
